\documentclass[11pt]{article}
\usepackage{geometry}
\geometry{a4paper, margin=1in}
\usepackage{setspace}
\usepackage{xcolor}
\usepackage{tikz}
\usepackage{hyperref}
\usepackage{enumitem}
\usepackage{needspace}
\usepackage{tabularx}
\hypersetup{
    colorlinks=true,
    linkcolor=blue,
    filecolor=magenta,
    urlcolor=cyan,
    pdftitle={Zhi Wang CV},
    pdfpagemode=FullScreen,
}
\linespread{0.99}
% Define custom symbols
\newcommand{\legendtriangle}[1][black]{\tikz \draw[#1,fill=#1] (0,0) -- (0.15,0.25) -- (0.3,0) -- cycle;}
\newcommand{\legendsquare}[1][black]{\tikz \draw[#1,fill=#1] (0,0) rectangle (0.25,0.25);}
\newcommand{\legendcircle}[1][black]{\tikz \draw[#1,fill=#1] (0,0) circle (0.125);}

\setlength{\parskip}{1em}
\usepackage{titlesec}
\titlespacing*{\section}{0pt}{0.06\baselineskip}{0.06\baselineskip}

\begin{document}

{\centering \textbf{\huge Zhi Wang} \par}
{\centering wangzhi0467@gmail.com $\vert$ \href{https://zhi0467.github.io}{Homepage} (zhi0467.github.io)$\vert$ \href{https://github.com/Zhi0467}{GitHub}\par}
\vspace{0.1em}

\section*{Education}
\hrule
\noindent
\textbf{University of California, San Diego} \hfill  San Diego, CA, USA\\
\textit{Ph.D. in Data Science} \hfill \textit{Sept. 2025 -- Ongoing}

\noindent
\textbf{University of Science and Technology of China} \hfill Hefei, China \\
\textit{B.S. in Mathematics, School of Gifted Young} \hfill \textit{Sept. 2021 -- Jun. 2025}

\noindent
\textbf{University of California, Berkeley} \hfill Berkeley, CA, USA \\
\textit{Visiting Undergraduate Student} \hfill \textit{Jan. 2023 -- Dec. 2023}

\vspace{0.2em}

\section*{Research Interests}
\hrule
\begin{itemize}[leftmargin=1.5em, topsep=2pt, partopsep=0pt, itemsep=0pt, parsep=0pt]
    \item \textbf{Reliable agentic systems:} monitoring, steering, and credit assignment for agents, and the infrastructure for scientists to collaborate with them.
    \item \textbf{Predictive science of optimization:} theory and forecasting of training dynamics, e.g. adaptive optimizers and predicting or monitoring learning metrics.
    \item \textbf{Embodied agents:} orchestrating reasoning and low-latency control for robotic tasks across different horizons.
\end{itemize}

\section*{Publications}
\hrule
\begin{itemize}
    \item Ke Zhang, Patricio Gallardo, Sudhir Murthy, Yi Xie, \textbf{Zhi Wang}, and Maziar Raissi, \textit{Beyond Compilation: Human-Calibrated Evaluation of Semantic Faithfulness in Lean Autoformalization}, NeurIPS 2026 MATH-AI Workshop. \href{https://openreview.net/forum?id=NyhNnJEJgL}{[Paper]}
    \item Eon Lee, Andrés R. Vindas-Meléndez, and \textbf{Zhi Wang}\footnote{In mathematics, authors are conventionally listed alphabetically, signifying equal contribution.}, \textit{Generalized Snake Posets, Order Polytopes, and Lattice-Point Enumeration}, Discrete Mathematics (2026): 115072.
\end{itemize}

\section*{Open Source Contributions}
\hrule
\noindent
\textbf{\href{https://github.com/Zhi0467/RCP}{RCP: Research Control Platform}}\vspace{-\parskip}
\begin{itemize}[leftmargin=1.5em, topsep=0pt, partopsep=0pt, itemsep=0pt, parsep=0pt]
    \item Harness and control plane over coding agents: one platform for research artifacts and asynchronous compute jobs, mediating how agents interact with humans and with compute.
    \item Conversational research assistant as the unified entry point: requests like ``start the experiment for me'' become managed asynchronous jobs, with human approval gating protected actions.
    \item Git-like research graph of questions, hypotheses, experiments, and evidence, with branching and merging.
    \item Runs as a macOS app or a self-hosted Linux team server with multi-member collaboration and attribution; agent lifecycles are verified by a TLA+-style state-space explorer.
\end{itemize}

\needspace{4\baselineskip}
\noindent
\textbf{Terminal-Bench}\vspace{-\parskip}
\begin{itemize}[leftmargin=1.5em, topsep=0pt, partopsep=0pt, itemsep=0pt, parsep=0pt]
    \item Contributed two benchmark tasks: \href{https://github.com/laude-institute/terminal-bench/pull/724}{merge-diff-arc-agi}, included in Terminal-Bench 2.0, and \href{https://github.com/laude-institute/terminal-bench/pull/936}{solve-sudoku}.
\end{itemize}

\section*{Current Research}
\hrule
\noindent
\textbf{Lean Prover Training: Whole-Proof Repair vs.\ Tree Search}\vspace{-\parskip}
\begin{itemize}[leftmargin=1.5em, topsep=0pt, partopsep=0pt, itemsep=0pt, parsep=0pt]
    \item Starting from the same checkpoint at equal compute, does a Lean prover learn more by writing whole proofs and repairing them from verifier feedback, or by proving step by step with tree search and expert iteration?
\end{itemize}

\noindent
\textbf{Retention vs.\ Plasticity in RL Post-Training}\vspace{-\parskip}
\begin{itemize}[leftmargin=1.5em, topsep=0pt, partopsep=0pt, itemsep=0pt, parsep=0pt]
    \item When a language model is trained with RL across domains, how do sequential, mixed-reward, consolidation-based and stabilized training trade off keeping skills in trained domains against staying able to learn new ones?
\end{itemize}

\noindent
\textbf{Predicting and Steering Chain-of-Thought Loops}\vspace{-\parskip}
\begin{itemize}[leftmargin=1.5em, topsep=0pt, partopsep=0pt, itemsep=0pt, parsep=0pt]
    \item Can a reasoning model's internal activations predict a chain-of-thought loop and steer the model away from it, and why do lightweight activation edits make models like Olmo-3 reason better on math?
\end{itemize}

\noindent
\textbf{Vision-Language Model Agents for Robot Manipulation}\vspace{-\parskip}
\begin{itemize}[leftmargin=1.5em, topsep=0pt, partopsep=0pt, itemsep=0pt, parsep=0pt]
    \item How can vision-language model agents perform robot manipulation faster and more cheaply, whether they control the robot directly or instruct a vision-language-action policy?
\end{itemize}

\noindent
\textbf{Self On-Policy Distillation for Multi-Turn Agents}\vspace{-\parskip}
\begin{itemize}[leftmargin=1.5em, topsep=0pt, partopsep=0pt, itemsep=0pt, parsep=0pt]
    \item Does extending self on-policy distillation to multi-turn agentic settings give a denser training signal that trains better small agents?
\end{itemize}


\section*{Prior Research Experience}
\hrule
\noindent\begin{tabularx}{\textwidth}{@{}p{5.5em}X@{}}
2024 -- 2025 & Simulating \textit{C. elegans} neural activity with connectome-constrained RNNs. \textit{Quan Wen, USTC.} \\[2pt]
2024 & Grokking and low-rank initialization in MLPs and transformers. \textit{Difan Zou, HKU.} \\[2pt]
2023 -- 2024 & Ehrhart theory of order polytopes; led to the \textit{Discrete Mathematics} paper. \textit{Andrés R. Vindas-Meléndez, UC Berkeley.} \\
\end{tabularx}

\section*{Teaching}
\hrule
\begin{itemize}[leftmargin=1.5em, topsep=2pt, partopsep=0pt, itemsep=0pt, parsep=0pt]
    \item Teaching Assistant, DSC 242: High-Dimensional Probability and Statistics, UC San Diego, Fall 2026.
    \item Teaching Assistant, DSC 140B: Representation Learning, UC San Diego, Winter 2026.
\end{itemize}

\end{document}
